The Gram test must apply to the actual span of allowed
operators. We verify the prefix bounds on that span before
using the chosen-plane projection in
Proposition~\ref{prop:residue-plane-neighbor-test}.
\paragraph{Prefix orders on the exact operator span.}
\label{par:plane-orbit}
\hypertarget{source-plane-orbit}{}
In the distinguished edge coordinates the flat generators of the model
are exactly \(B_b\,\mathrm{SPR}\). Indeed pulling back each row of the
true frame determinant by the edge substitution gives its designated
plane row, so the product \(\mathrm{SPtrue}\) pulls back to
\(\mathrm{SPR}\). Every monomial in this common factor uses only
even plane coordinates. We work over the fraction field in their
linear span and apply words in \(X_{sa}\partial_{sb}\), using each
top \(s\) at most \(c_0\) times. Conjugation by an invertible linear
substitution within each letter takes a row operator to a linear
combination of row operators at the same letter. In matrix-unit notation the substitution by \(g\) acts on a row operator through
\[
 gE_{ab}g^{-1}=\sum_{c,d}g_{ca}(g^{-1})_{bd}E_{cd}.
\]
Each term acts at the same top, so the word allowance persists in edge coordinates.

List the tops after the incoming frame of segment \(j\) in increasing
order as \(x_0,\ldots,x_{h+r-1}\), where \(h\) is the number strictly
between the frames. The last \(r\) positions are the outgoing frame,
with shared index
\[
 m=r-1-(u-h)
\]
at position \(x_u\).
Write \(B_u\) for the true degree of the oldest of the \(r\) even
predecessors of \(x_u\), divided by its radix in this segment. Put
\[
 s_u=\lfloor B_u/(Yk^2)\rfloor.
\]

For a joint complement type,
let \(k_u\) denote its largest weight at the prefix of the first
\(u\) of the listed tops. In this paragraph frame subscripts on
\(B_u,s_u,k_u\) are implicit. Define \(T_u\) as the largest of
the following applicable quantities:
\[
 \frac{B_u}{2},\quad k_{u+1}+m_*+s_u\quad(x_u\text{ even}),\quad
 k_{u-1}+m_*+s_u\quad(u>0,\ x_{u-1}\text{ odd}).
\]
The chain and link thresholds for the plane Gram argument are
\begin{equation}\label{eq:plane-chain-link-thresholds}
   T_u-s_u/8\quad(0\le u\le h),\qquad
   T_u-3s_u/8\quad(h<u<h+r),
\end{equation}
respectively. Its separate chain
slack is \(3s_u/8\). We show that these thresholds work for
every choice of low digits, quotient multidegrees, joint
complement type and coefficient of that type.

\begin{proposition}[Prefix thresholds on the full permitted operator span]\label{prop:plane-operator-prefix-orders}
On the fraction-field span of the flat generators \(B_b\,\mathrm{SPR}\)
and the allowed words in \(X_{sa}\partial_{sb}\) with at most
\(c_0\) operators at each top, the chain and link thresholds from \eqref{eq:plane-chain-link-thresholds} are
\[
 \begin{cases}T_u-s_u/8,&0\le u\le h,\\T_u-3s_u/8,&h<u<h+r.\end{cases}
\]
 They hold after taking any indicated digit class,
quotient multidegree and joint complement type, with the same
allowance \(s_u\) for the selected coefficient.
\end{proposition}

\begin{proof}
\proofstep{Expansion into segment generators.}
First expand a generator in its presentation by products of
segment pairs \(A_jG_j\) in the plane and complement
coordinates. Do this before selecting a coefficient. The external
label \(l\) at each top is strictly below the segment radix there.
Take homogeneous components at the prescribed true degrees of
\(B_b\) in this expansion. Renamed products of homogeneous segment
factors have definite true degrees, and the sum of the components
in other degrees is zero. Thus the degree bound below applies on
each retained product. Expand \(\mathrm{SPR}\) into monomials
\(X^\gamma\).
In a coordinate of radix \(P\) and prescribed residue \(\rho\), the term can survive extraction only if
\[
 l+\gamma\equiv\rho\pmod P,\qquad (l+\gamma-\rho)/P\in\mathbb Z_{\ge0}.
\]
Apply this test in every coordinate.
In the matching class the quotient after multidegree projection
is a tensor product. Its complement factor is the projected
product of the \(G_j\) with their positions renamed.

The
complement is absent on incoming rows and the complement label
has zero quotient. Its plane factor, before multidegree projection,
is
\[
        X^b\,\phi\!\left(\prod_j \chi_j A_j\right).
\]
Here \(b\) is the plane quotient of \(l+\gamma\) by the low
radices, coordinate by coordinate. It vanishes on odd letters.
The variable map \(\chi_j\) sends intervening, outgoing, and
incoming letters respectively to the block letters,
\(L_{j,m}\), and \(W_{j-1,m}\). If \(P_{\rm low}\) is the low
radix on the preceding frame, the exponent multiplier at an incoming
row is
\[
           P_{\rm coarse}=P_{\rm low}\,P_{j-1,m},
                 \qquad P_{j-1,m}=p^{z_j-z_{j-1}}.
\]
This proves the displayed formula after collapse, since \(\phi\)
raises \(W_{j-1,m}\) to the \(P_{j-1,m}\)-th power.
These maps use the original frame radices after the transport in
Section~\ref{par:plane-rank}.

If the complement projection vanishes, the extracted plane
coefficient is zero. Otherwise its type at every prefix has
largest weight at most the upper Laurent order of the corresponding
prefix of \(G_j\). The other complement factors are outside the
cut, and the multidegree projection commutes with its action.
An isotypic projection with larger first weight would annihilate
the vector by the conjugate torus criterion. The original segment
inequalities now give
\[
 \operatorname{ord}_{<x_u} A_j\ge T_u.
\]

For an even letter use the complement through that letter. For a
preceding odd letter use its complement before it and the plane
through it. On segment support this last plane prefix is precisely
the prefix before \(x_u\). Its oldest even predecessor and hence
the available margin are unchanged in traversing the odd letter.
The absolute plane lower bound gives the first term of the maximum.
The extracted coefficient is a scalar multiple of the projected
plane factor. All these conditions describe linear spaces.

\proofstep{Transport of the prefix cuts by relabeling.}
For the chain cut \(S_{ju}\), the renaming \(\chi_j\) carries the
prefix before \(x_u\) to all incoming \(W\)'s and the first \(u\)
intervening letters with their original signs. Other segments
have no variables in this cut. Thus
\[
 \Pi=\prod_j\chi_j A_j
\]

has order at least \(T_u\). Restrict to a desired true quotient
multidegree of \(X^b\phi\Pi\) before applying the shared digit
extraction. On \(\Pi\) this restriction selects the model
multidegrees for which the weighted sum of letter degrees after
collapse has the prescribed value. Their projections commute with
the chain action. Expand the resulting polynomial by monomials
\(L^M\) of the label letters, which are outside the cut. Each
coefficient, a polynomial in block and \(W\) letters, retains the
order.

For each shared coordinate \(t\), divide the exponent in
\(X^b\phi(L^M)\) by \(P_t\):
\[
     b_t+M_t=\rho_t+P_t a_t,\qquad 0\le\rho_t<P_t,\quad a_t\ge0 .
\]
Thus the coefficient contributes only to the residue \(\rho_t\)
and acquires the fixed carry \(W_t^{a_t}\). For a tested residue \(\rho'_t\), the carry factor after extraction is therefore
\[
 \mathbf1_{\rho'_t=\rho_t}W_t^{a_t}.
\]
Apart from these carries the extra monomial
is the block part of \(X^b\). All affected coordinates of this monomial are vectors.
They have nonnegative least weight and multiplication preserves
the lower bound. This proves the required chain order for every
shared residue.

A true link at \(x_u\) with \(u>h\) sees the intervening letters
and the earlier outgoing rows. The original prefix also sees the
incoming rows. In a homogeneous component of \(A_j\), let their
total degree be \(d_W\). Let \(D(T)\) be the rank-one one-parameter torus
used for the original Laurent substitution.

To omit the
incoming variables \(W\) from this cut, substitute \(D(T)^{-1}W\)
in the transformed polynomial. A monomial of degree \(d_W\) accumulates \(d_W\) inverse weights, each either zero or \(-1\). Its Laurent order in this substitution is therefore
\[
 -\sum_{a=1}^{d_W}\varepsilon_a\ge-d_W,\qquad \varepsilon_a\in\{0,1\}.
\]
Therefore
\[
 \operatorname{ord}_{\rm prefix\ without\ incoming}(A_j)
 \ge \operatorname{ord}_{\rm original\ prefix}(A_j)-d_W .
\]
Coefficient extraction in the now untransformed \(W\)
preserves this bound for each remaining polynomial. Under collapse those incoming monomials go to frame
\(j-1\), outside this true link. Here \(J\ge2\) follows from the
frame span bound \(G'\) and \(3G'\le n'\), since one frame would
recur only after \(n'\) positions. All factors from other
segments and from \(b\) that meet this link use vectors there.
Projection to the required multidegree preserves the lower bound.

Only an incoming degree \(d_m^W\) with
\[
 P_{j-1,m}d_m^W\le (d_{\rm pl}(j-1,m)-|b_{j-1,m}|)_+
\]

can contribute. At the row \(y_m\) of that preceding frame put
\(t_m=z_{j-1}-y_m\), \(g=z_j-z_{j-1}\). Removing the quotient
of the external monomial leaves at most the original generator
degree \(2a\) at true scale. Indeed after multiplying the
remaining quotient by \(p^{t_m}\) the degree contributed by
the segment products is bounded by their prescribed true degree,
which is at most the flat degree. The \(t_m\) are distinct and
nonnegative. Consequently
\[
 d_W\le\sum_m 2a p^{-g-t_m}\le 3a p^{-g}.
\]

If \(y\) is the oldest even predecessor of the current outgoing
row, there are at least \(r\) intervening even tops before it is
tested. The first intervening top is more than \(L_1\) after
\(z_{j-1}\). Hence \(y\ge z_{j-1}+L_1+1\), and
\(a\le B_u p^{z_j-y}\). The host spacing gives
\(p^{L_1+1}\ge48Yk^2\); the size bound gives
\(B_u/(Yk^2)\ge1\), since \(B_u\ge a/Q\).
The permitted incoming degree is therefore bounded by
\[
 d_W\le \frac{3B_u}{p^{L_1+1}}
     \le \frac{B_u}{16Yk^2}\le s_u/8 .
\]
The last step uses \(B_u/(Yk^2)\le2s_u\).
Subtracting this degree from the full prefix order \(T_u\)
gives the initial link bound.

\proofstep{Admitting row operators in the order estimate.}
It remains to allow row operators. For one coordinate of low radix
\(P\), write its true power as \(v+Pq\), \(0\le v<P\).
Write \(x^v F(X)\) with \(X=x^P\). The two derivative
cases follow from the chain rule:
\[
 \partial_x\!\bigl(x^vF(x^P)\bigr)=
 \begin{cases}
   x^{v-1}(v+P X\partial_X)F(X),&v>0,\\
   x^{P-1}P\partial_X F(X),&v=0 .
 \end{cases}
\]
For multiplication the two cases are
\[
 x\,x^v F(X)=
 \begin{cases}
 x^{v+1}F(X),&v+1<P,\\
 X F(X),&v=P-1 .
 \end{cases}
\]
At \(P=1\) the second case applies to \(v=0\). The same
formulas at the second digit extraction act on the incoming
\(W\) quotient, or on the block coordinate at radix one. On an exponent \(v+Pq\), the Euler operator \(x\partial_x\) acts by \(v+Pq\). Thus on the quotient coordinate \(y\) in the residue class \(v\) its action is
\[
 v+P\,y\partial_y.
\]

Thus each multiplication or differentiation before extraction costs
at most one in a lower cut order that meets the letter. Under
conjugation the torus on a coordinate has weights \(0,1\) on a
vector and \(0,-1\) on a covector, with opposite signs on the
derivatives. The Euler operator also costs at most one, since one
of its two factors has nonnegative least weight. On a disjoint cut
these operations preserve the order.

For a complement operation, insert the joint decomposition of the
input before the operation and project the output. At each prefix \(I\), the upper-order estimate for an
operator at \(x\) allows a transition from \(k_I\) only to
\[
                          k'_I\le k_I+\mathbf1_{x\in I}.
\]
Complement coefficients commute with the plane cuts. Thus
this transition changes the maximum defining \(T_u\) by
at most one for a complement operation. It can
occur only for a top at or before \(x_u\) in that segment.
A plane operation can affect the shared \(W\) cut even if its
letter belongs to the previous frame, but that row occurs in the
incoming part of the segment, also before \(x_u\).
This accounts for both occurrences of a shared row, including
across the cyclic endpoint.

\proofstep{The exact operator span.}
If \(a_u\) operators in the word act at such tops up to the tested position \(u\), induction therefore bounds their cumulative loss by 
\[
 \underbrace{2+\cdots+2}_{a_u\,\text{operators}}=2a_u.
\]
 The
cut and type estimates for a single half-operation apply in all
input digit classes and degrees. The induction therefore survives
both quotient-degree changes and the sum over input types. For one cut, at most \(N_1c_0\) operators can touch its letters.
The host inequalities give
\[
 a\ge16N_1c_0Yk^2Q,
\]
and
\[
 B_u\ge a/Q
\]
implies
\[
 s_u\ge16N_1c_0.
\]

For a word let \(m_I\) be the number of operators meeting
the cut. The induction allows at most \(2m_I\) in the
chain and at most \(4m_I\) in the link, with both occurrences
of a shared row included. The budget \(m_I\le N_1c_0\) and
\(16N_1c_0\le s_u\) in particular give
\[
    \operatorname{ord}_{\rm chain}\ge T_u-s_u/8,\qquad
    \operatorname{ord}_{\rm link}\ge T_u-s_u/8-s_u/4 .
\]
The first loss is charged from the starting chain threshold
and the second from the initial link threshold.
These link thresholds are positive since \(T_u\ge B_u/2\)
and \(s_u\le B_u/(Yk^2)\). Linear combinations over the fraction field
preserve these inequalities. This proves the membership used for
the selected plane polynomial in Section~\ref{par:plane-flat-test}.
 \label{cor:plane-orbit-1}
\end{proof}
